\section*{Capítulo \ref{ch:b1:plate-kinematics} --- Cinemática de placas sobre una esfera}

\begin{solution}{exo:b1:plate-kinematics:1}
(a) $\qty{0.210}{\degree\per\Myr} = 0.210\times\pi/180 = \qty{3.665e-3}{rad}$
por millón de años, es decir \qty{3.665e-9}{rad.yr^{-1}}; con
$\qty{1}{yr} = \qty{3.156e7}{s}$, \qty{1.161e-16}{rad.s^{-1}}. (b) $360/0.210 =
\qty{1714}{\Myr}$, más de un tercio de la edad de la Tierra. (c) $\omega R =
\num{3.665e-9}\times\qty{6.371e9}{mm} = \qty{23.35}{mm.yr^{-1}}$.
\end{solution}

\begin{solution}{exo:b1:plate-kinematics:2}
$v = \omega R\sin\Delta = 23.35\sin\Delta$: \qty{11.7}{mm.yr^{-1}} a \ang{30},
\qty{20.2}{mm.yr^{-1}} a \ang{60} y \qty{23.35}{mm.yr^{-1}} a \ang{90}. La
expansión es más rápida donde la frontera está más lejos del polo, hasta
\ang{90}: entre Eurasia y Norteamérica, en el extremo sur de la frontera, cerca de las Azores, a \ang{78} del polo.
\end{solution}

\begin{solution}{exo:b1:plate-kinematics:3}
El subcron dura $1.072 - 0.988 = \qty{0.084}{\Myr}$, y una \omterm{def:b1:plate-kinematics:spreading-rate}{velocidad de expansión media} en
\unit{mm.yr^{-1}} es también una velocidad en \unit{km\per\Myr}. (a) $50\times0.084 =
\qty{4.2}{km}$ a \qty{50}{mm.yr^{-1}}; $10\times0.084 = \qty{0.84}{km}$ a
\qty{10}{mm.yr^{-1}}. (b) La franja se reparte a ambos lados del eje:
$2\times10\times0.781 = \qty{15.6}{km}$.
\end{solution}

\begin{solution}{exo:b1:plate-kinematics:4}
$66/2.581 = \qty{25.6}{mm.yr^{-1}}$ hacia el oeste y
$64/2.581 = \qty{24.8}{mm.yr^{-1}}$ hacia el este; la velocidad total es
$130/2.581 = \qty{50.4}{mm.yr^{-1}}$. La expansión es casi simétrica: el flanco oeste
ha crecido un \qty{3}{\percent} más rápido.
\end{solution}

\begin{solution}{exo:b1:plate-kinematics:5}
$\cos\Delta = \sin\ang{61.8}\sin\ang{85} + \cos\ang{61.8}\cos\ang{85}
\cos(\ang{0} - \ang{139.6}) = 0.8813\times0.9962 + 0.4726\times0.08716\times
(-0.7615) = 0.87795 - 0.03137 = 0.84658$, así que $\Delta = \ang{32.16}$ y
$v = 23.35\sin\ang{32.16} = \qty{12.4}{mm.yr^{-1}}$. Esto está un poco por encima de los
\qtyrange{11}{12}{mm.yr^{-1}} medidos, pero el \omterm{def:b1:plate-kinematics:euler-pole}{vector de Euler} es un único mejor ajuste a
453 velocidades, de las que las velocidades se apartan \qty{0.7}{mm.yr^{-1}} de media, y las velocidades individuales cerca de la
lenta dorsal de Gakkel están entre las menos precisas.
\end{solution}

\begin{solution}{exo:b1:plate-kinematics:6}
(a) $\cos\Delta = \sin\ang{61.8}\sin\ang{39.44} + \cos\ang{61.8}\cos\ang{39.44}
\cos(\ang{169.43}) = 0.20111$: $\Delta = \ang{78.40}$. (b) $v = 23.35
\sin\ang{78.40} = \qty{22.87}{mm.yr^{-1}}$, el \qty{23}{mm.yr^{-1}} medido
justo al norte de la unión. (c) Con la fórmula del método, el numerador es
$\sin\ang{169.43}\cos\ang{61.8} = 0.08668$ y el denominador
$\cos\ang{39.44}\sin\ang{61.8} - \sin\ang{39.44}\cos\ang{61.8}\cos\ang{169.43}
= 0.97573$: el polo se ve con $\beta = \ang{5.08}$, casi exactamente al norte, y
Eurasia se mueve hacia $\ang{95.1}$. (d) Las transformantes son paralelas al
movimiento relativo: se orientan hacia \ang{095}, casi este--oeste.
\end{solution}

\begin{solution}{exo:b1:plate-kinematics:7}
(a) La velocidad angular del \omterm{def:b1:plate-kinematics:euler-pole}{polo de Euler} es $\vec\omega = \omega(\cos\varphi_p\cos\lambda_p, \cos\varphi_p\sin\lambda_p,
\sin\varphi_p)$: $\vec\omega_{\mathrm{EU\,NA}} = (-0.0756, 0.0643, 0.1851)$ y
$\vec\omega_{\mathrm{NB\,NA}} = (0.0333, 0.0282, 0.2289)$, en
\unit{\degree\per\Myr}. (b) $\vec\omega_{\mathrm{NB\,EU}} =
\vec\omega_{\mathrm{NB\,NA}} - \vec\omega_{\mathrm{EU\,NA}} = (0.1089, -0.0361,
0.0438)$, de longitud
\[
  \sqrt{0.011859 + 0.001303 + 0.001918} = \qty{0.1228}{\degree\per\Myr}.
\]
Su polo está en la \omterm{def:b1:spherical-astronomy:geographic}{latitud}
$\arcsin(0.0438/0.1228) = \ang{20.9}$\,N y la longitud
$\arctan(-0.0361/0.1089) = \ang{-18.3}$, es decir \ang{18.3}\,W (la
componente $x$ es positiva, así que no se suma \ang{180}). (c) La frontera entre
Nubia y Eurasia ofrece pocas medidas directas: entre las Azores y
Gibraltar, solo los \omterm{def:b1:spherical-astronomy:horizontal}{acimutes} de la falla de Gloria. Las dos expansiones
a través del Atlántico están medidas por centenares de velocidades, y su diferencia
hereda su precisión.
\end{solution}

\begin{solution}{exo:b1:plate-kinematics:8}
(a) Alrededor de $\vec e_z$ en \ang{90}: $\cos\ang{90} = 0$, $\vec e_z\cdot\vec e_x = 0$
y $\vec e_z\times\vec e_x = \vec e_y$, así que $\vec e_x$ va a $\vec e_y$. Luego
alrededor de $\vec e_x$ en \ang{90}: $\vec e_x\cdot\vec e_y = 0$ y
$\vec e_x\times\vec e_y = \vec e_z$, así que $\vec e_y$ va a $\vec e_z$: el punto
termina en el polo norte. (b) Primero alrededor de $\vec e_x$: $\vec e_x$ está sobre el eje,
$\vec e_x\times\vec e_x = \vec 0$ y $(1 - \cos\ang{90})(\vec e_x\cdot\vec e_x)
\vec e_x = \vec e_x$, así que no se mueve; luego, alrededor de $\vec e_z$, va a
$\vec e_y$, a \ang{0}\,N, \ang{90}\,E. (c) Los resultados difieren: las \omterm{def:b1:plate-kinematics:finite-rotation}{rotaciones finitas} no conmutan. Las reconstrucciones sucesivas deben componerse
multiplicando las matrices de rotación en el orden correcto, nunca sumando ángulos
o polos.
\end{solution}

\begin{solution}{exo:b1:plate-kinematics:9}
(a) En el espacio de velocidades ponga $a$ en el origen, $b$ a \qty{30}{mm.yr^{-1}} hacia
el este y $c$ a \qty{40}{mm.yr^{-1}} hacia el sur. (b) $\vec v_{CB} =
\vec v_{CA} - \vec v_{BA} = (0, -40) - (30, 0) = (-30, -40)$ (componentes este y norte):
\qty{50}{mm.yr^{-1}} hacia el \omterm{def:b1:spherical-astronomy:horizontal}{acimut}
$\arctan\bigl((-30)/(-40)\bigr) + \ang{180} = \ang{216.9}$, suroeste. (c) La
dorsal es perpendicular a esta dirección: su orientación es
$\ang{216.9} - \ang{90} = \ang{126.9}$, es decir noroeste--sureste.
\end{solution}

\begin{solution}{exo:b1:plate-kinematics:10}
(a) En la matriz de Rodrigues, la identidad aporta $3\cos\alpha$, la matriz
$\vec n\,\vec n^{\mathsf T}$ la traza $n_x^2 + n_y^2 + n_z^2 = 1$, multiplicada por
$1 - \cos\alpha$, y la matriz antisimétrica nada: $\operatorname{tr}R =
3\cos\alpha + 1 - \cos\alpha = 1 + 2\cos\alpha$. (b) Aquí $\operatorname{tr}R = 0$,
así que $\cos\alpha = -1/2$ y $\alpha = \ang{120}$. El eje cumple $R\vec n =
\vec n$: $R(x, y, z) = (z, x, y)$, así que $x = y = z$ y $\vec n =
(1, 1, 1)/\sqrt3$, en \omterm{def:b1:spherical-astronomy:geographic}{latitud} $\arcsin(1/\sqrt3) = \ang{35.26}$\,N y
\omterm{def:b1:spherical-astronomy:geographic}{longitud} \ang{45}\,E. (c) $R\vec e_x = \vec e_y$, $R\vec e_y = \vec e_z$,
$R\vec e_z = \vec e_x$: los tres ejes se permutan, como lo hace un tercio de vuelta alrededor de la diagonal. La fórmula de Rodrigues con $\alpha = +\ang{120}$ alrededor de $\vec n$
en efecto da $\vec e_x\mapsto -\tfrac12\vec e_x + \tfrac12(1, 1, 1) +
\tfrac{\sqrt3}{2}\cdot\tfrac{1}{\sqrt3}(0, 1, -1) = (0, 1, 0)$, así que la rotación
es antihoraria vista desde encima de \ang{35.26}\,N, \ang{45}\,E.
\end{solution}

\begin{solution}{exo:b1:plate-kinematics:11}
(a) En una transformante la velocidad relativa $\vec\omega\times\vec r$ va a lo largo de la
falla, luego paralela a $\vec t$; es perpendicular a $\vec\omega$, de donde
$\vec t\cdot\vec\omega = 0$. El plano perpendicular a $\vec t$ que pasa por el
centro contiene $\vec r$ (ya que $\vec t$ es horizontal) y $\vec\omega$: es
el plano del \omterm{def:b1:spherical-astronomy:spherical-triangle}{círculo máximo} que pasa por el punto y es perpendicular a la falla,
y el polo yace sobre ese círculo. Dos planos así se cortan a lo largo de la recta de
$\vec t_1\times\vec t_2$. (b) $\vec t_1 = (0.5436, 0.8367, -0.0673)$ en la
unión y $\vec t_2 = (0.4942, 0.8639, -0.0968)$ en Islandia, así que
$\vec t_1\times\vec t_2 = (-0.02285, 0.01936, 0.05612)$, de longitud $0.06362$.
Su \omterm{def:b1:spherical-astronomy:geographic}{latitud} es $\arcsin(0.05612/0.06362) = \ang{61.9}$\,N y su longitud
$\ang{180} - \arctan(0.01936/0.02285) = \ang{139.7}$\,E (componente $x$ negativa):
el polo del capítulo. (c) Mal: los dos \omterm{def:b1:spherical-astronomy:spherical-triangle}{círculos máximos} se cortan con un ángulo pequeño,
ya que las orientaciones difieren solo en \ang{8} a lo largo de \ang{26} de frontera.
Cambiar cualquiera de las orientaciones en un grado desplaza el polo unos \ang{15}. Las determinaciones reales combinan muchas transformantes con velocidades de expansión,
cuya variación a lo largo de la frontera fija la distancia angular al polo.
\end{solution}

\begin{solution}{exo:b1:plate-kinematics:12}
(a) $(94 - 4)/3 = \qty{30}{mm.yr^{-1}}$. (b) Las diferencias separadas exactamente un año son
$24 - 4 = 20$, $25 - 5 = 20$, $26 - 6 = 20$, $65 - 15 = 50$, $74 - 24 = 50$,
$75 - 25 = 50$, $76 - 26 = 50$, $85 - 65 = 20$ y $94 - 74 = 20$ (mm). Ordenadas:
cinco valores de 20 y cuatro de 50; la mediana es \qty{20}{mm.yr^{-1}}. (c) Dos
posiciones separadas un año están en la misma fase de la oscilación estacional,
que se cancela en su diferencia. Un salto afecta solo a los pares que lo cruzan;
aquí, cuatro de los nueve. Mientras menos de la mitad de los pares cruce un salto,
la mediana lo ignora, mientras que la estimación de los extremos del apartado (a) toma todo el salto
como movimiento y resulta un \qty{50}{\percent} demasiado alta. El método falla cuando
los saltos son tan frecuentes que la mayoría de los pares anuales contienen uno, o cuando la estación
se mueve de forma irregular, por ejemplo tras un gran terremoto.
\end{solution}

\begin{solution}{pb:b1:plate-kinematics:1}
\textbf{1.} $\omega = 0.210\times(\pi/180)/\num{e6} = \qty{3.665e-9}{rad.yr^{-1}}$;
$\omega R = \num{3.665e-9}\times\qty{6.371e9}{mm} = \qty{23.35}{mm.yr^{-1}}$.

\textbf{2.} $\cos\Delta = \sin\ang{61.8}\sin\ang{64.5} + \cos\ang{61.8}
\cos\ang{64.5}\cos(\ang{-157.6}) = 0.79546 - 0.18809 = 0.60737$, así que
$\Delta = \ang{52.60}$.

\textbf{3.} $v = 23.35\sin\ang{52.60} = 23.35\times0.7944 =
\qty{18.55}{mm.yr^{-1}}$.

\textbf{4.} El numerador es $\sin\ang{157.6}\cos\ang{61.8} = 0.18011$ y el
denominador
\[
  \cos\ang{64.5}\sin\ang{61.8} - \sin\ang{64.5}\cos\ang{61.8}\cos\ang{157.6} =
  0.37940 + 0.39436 = 0.77376;
\]
ambos son positivos, así que $\beta = \arctan 0.2328 =
\ang{13.1}$. Eurasia se mueve hacia $\ang{13.1} + \ang{90} = \ang{103.1}$.

\textbf{5.} Este: $18.55\sin\ang{103.1} = \qty{18.07}{mm.yr^{-1}}$; norte:
$18.55\cos\ang{103.1} = \qty{-4.20}{mm.yr^{-1}}$, es decir \qty{4.20}{mm.yr^{-1}}
hacia el sur.

\textbf{6.} $\cos\Delta = 0.8851$, $\Delta = \ang{27.74}$, $v = 23.35
\sin\ang{27.74} = \qty{10.9}{mm.yr^{-1}}$. Las dos placas giran alrededor de un polo en
el este de Siberia: la dorsal de Gakkel está cerca de él, y la velocidad cae con
$\sin\Delta$ aunque las placas sean las mismas.

\textbf{7.} Oeste: $94/9.779 = \qty{9.61}{mm.yr^{-1}}$; este: $97/9.779 =
\qty{9.92}{mm.yr^{-1}}$.

\textbf{8.} $(94 + 97)/9.779 = 191/9.779 = \qty{19.53}{mm.yr^{-1}}$.

\textbf{9.} Desde C2An: $54.1/2.581 = \qty{20.96}{mm.yr^{-1}}$; desde la base de
C1n: $19.8/0.781 = \qty{25.35}{mm.yr^{-1}}$. No concuerdan: cuanto más joven es el límite, mayor es la velocidad. Un cambio real de
velocidad es posible, pero un exceso constante en las distancias, el mismo a todas las edades, produciría exactamente esto.

\textbf{10.} La razón entre las distancias del oeste y del este es de 0.96 a 0.98: el
flanco este ha crecido algunos por ciento más. La velocidad total depende solo de la
suma de las dos distancias, así que la asimetría no la cambia.

\textbf{11.} $(191 - 5.0)/9.779 = \qty{19.02}{mm.yr^{-1}}$;
$(54.1 - 5.0)/2.581 = \qty{19.02}{mm.yr^{-1}}$; $(19.8 - 5.0)/0.781 =
\qty{18.95}{mm.yr^{-1}}$. Los tres concuerdan ahora dentro de
\qty{0.1}{mm.yr^{-1}}. Los mismos \qty{5}{km} son un cuarto de los \qty{19.8}{km}
del par más joven, pero menos del \qty{3}{\percent} de los \qty{191}{km} del más antiguo;
los límites antiguos dan las velocidades más fiables.

\textbf{12.} La velocidad media corregida es $186/9.779/2 = \qty{9.51}{mm.yr^{-1}}$,
y $9.51\times0.084 = \qty{0.80}{km}$, un tercio de los \qty{2.5}{km} en que
se desplaza cada límite. Se esperaría que una franja tan estrecha se difuminara.
Que se reconozca en la dorsal de Reykjanes, mientras que falta en la dorsal de
Gakkel, escasa de magma, sugiere que un suministro abundante de magma registra el campo
fielmente incluso en franjas estrechas.

\textbf{13.} Las velocidades corregidas, de 18.95 a \qty{19.02}{mm.yr^{-1}}, concuerdan con
los \qty{19.0}{mm.yr^{-1}} del vector de placa; sin la corrección lo superaban en
un 3 a un 33 por ciento.

\textbf{14.} $\cos\Delta = \sin\ang{64.139}\sin\ang{59.018} +
\cos\ang{64.139}\cos\ang{59.018}\cos\ang{27.554} = 0.97054$, así que $\Delta =
\ang{13.94}$; la recta es $2R\sin(\Delta/2) = 12742\times\sin\ang{6.97} =
\qty{1546}{km}$.

\textbf{15.} $1000/17.5 = 57$ años.

\textbf{16.} Las dos direcciones difieren en $\ang{123.1} - \ang{114.8} =
\ang{8.3}$: la componente horizontal a lo largo del \omterm{def:b1:spherical-astronomy:spherical-triangle}{círculo máximo} es
$18.9\cos\ang{8.3} = \qty{18.70}{mm.yr^{-1}}$, y a lo largo de la recta
$18.70\cos\ang{6.97} = \qty{18.56}{mm.yr^{-1}}$, unos \qty{18.6}{mm.yr^{-1}}.

\textbf{17.} La recta se alarga \qty{17.5}{mm.yr^{-1}}, \qty{6}{\percent} menos que los \qty{18.6}{mm.yr^{-1}} predichos. Ambas
estaciones están lejos de la frontera de placas, así que la diferencia no es local. O bien las placas se mueven hoy algo más despacio
que de media en los últimos tres millones de años, o una de las dos estimaciones lleva un pequeño error sistemático; las medidas por
sí solas no lo dicen. Un \omterm{def:b1:plate-kinematics:euler-pole}{vector de Euler} calculado solo con posicionamiento por satélite predice \qty{17.1}{mm.yr^{-1}} para esta
recta, cerca del valor medido.

\textbf{18.} Este: $21.12 - (-1.07) = \qty{22.19}{mm.yr^{-1}}$; norte:
$-6.97 - (-0.97) = \qty{-6.00}{mm.yr^{-1}}$. Velocidad $\sqrt{22.19^2 + 6.00^2} =
\qty{23.0}{mm.yr^{-1}}$, hacia $\ang{180} - \arctan(22.19/6.00) =
\ang{105.1}$.

\textbf{19.} La dirección es cercana a los \ang{103} del vector de placa, pero la velocidad supera sus \qty{18.6}{mm.yr^{-1}} en un
cuarto. Ambas estaciones están dentro de la zona donde la frontera se deforma, y Höfn, cerca del mayor casquete, se mueve con el
rebote de la \omterm{def:b1:layered-earth:layers}{corteza} a medida que el hielo adelgaza: hacia arriba \qty{12.9}{mm.yr^{-1}}, y lateralmente, alejándose del hielo. Las
estaciones dentro de la zona de frontera miden el movimiento de la placa más la deformación local.

\textbf{20.} $0.210\times10 = \ang{2.10}$.

\textbf{21.} El punto recorre un \omterm{def:b1:spherical-astronomy:spherical-triangle}{círculo menor} de radio $R\sin\Delta$:
$6371\times0.7944\times(2.10\times\pi/180) = \qty{185.5}{km}$. La anchura corregida de
la \omterm{def:b1:layered-earth:layers}{corteza} formada desde el borde joven de C5n, \qty{186}{km} en
\qty{9.8}{Myr}, es casi la misma: el movimiento ha sido estable durante diez millones de años.

\textbf{22.} Componentes (este, norte): Eurasia
$(22.9\sin\ang{95.1}, 22.9\cos\ang{95.1}) = (22.81, -2.04)$, Nubia
$(19.2\sin\ang{103.8}, 19.2\cos\ang{103.8}) = (18.65, -4.58)$. Eurasia respecto
de Nubia: $(4.16, 2.54)$, es decir \qty{4.9}{mm.yr^{-1}} hacia
$\arctan(4.16/2.54) = \ang{58.6}$, este-nordeste.

\textbf{23.} Una unión de tres dorsales es estable sean cuales sean las velocidades. La
velocidad de la pregunta 22 hace que Eurasia y Nubia se separen, de este-nordeste a
oeste-suroeste: su frontera en la unión debe abrirse, lentamente; es el rift de Terceira,
cuya orientación este-sureste hace que la apertura sea oblicua.

\textbf{24.} Parte I: \qty{18.6}{mm.yr^{-1}} en el centro de Islandia
(\qty{19.0}{mm.yr^{-1}} en el perfil de la Parte II). Parte II: 18.95 a
\qty{19.02}{mm.yr^{-1}} una vez corregidas. Parte III: la recta Reikiavik--Stavanger
se alarga \qty{17.5}{mm.yr^{-1}}, donde el vector de placa predice
\qty{18.6}{mm.yr^{-1}}; el par dentro de Islandia da \qty{23.0}{mm.yr^{-1}},
que no es el movimiento de la placa. El vector de placa y las franjas concuerdan; los
satélites difieren en un \qty{6}{\percent}. Las franjas promedian hasta unos diez
millones de años, el vector de placa sobre \qty{3.16}{Myr}; los satélites miden el
movimiento de los últimos veinte a treinta años.

\textbf{25.} El vector de placa da \qty{18.55}{mm.yr^{-1}} a través de Islandia, y las
franjas corregidas dan \qty{19.0}{mm.yr^{-1}} junto a él: con precisión de milímetro por año,
la \omterm{def:b1:plate-kinematics:spreading-rate}{velocidad de expansión total} a través de Islandia es
\textbf{\qty[text-series-to-math=true]{19}{mm.yr^{-1}}}.
\end{solution}
